\section{Hypothesis Test}

% \paragraph{Data and design.}
To test whether 311 adoption is associated with changes in city government spending,
adoption dates from Section~1 are combined with the Census Annual Survey of
Government Finances, Individual Unit Data, for 2017 through 2022
\citep{census_govfin2022}.
Total general expenditure per capita is constructed from the item-level fixed-width
files by summing current-operations (E-prefix), capital-construction (F-prefix), and
other-capital-outlay (G-prefix) items, plus general interest on debt (I89), while
excluding utility enterprise operations.
Cities are matched to the finance data directly via Census place FIPS code,
achieving a 100\,\% match for the 153-city regression sample (99 adopted with known
year plus 54 never-adopted cities).
The estimator is two-way fixed effects (city and year fixed effects), with standard
errors clustered at the city level.
The primary data constraint is panel length: Individual Unit Finance files are
publicly available only for 2017 through 2022, so cities that adopted before 2017
appear as always-treated throughout the panel with no observable pre-period.
The 148 cities with confirmed 311 but unknown adoption timing are excluded from
the regression because the post-treatment indicator requires a specific year.

% \paragraph{Baseline result.}
Table~\ref{tab:twfe} reports estimates from two specifications.
In the full sample (867 city-year observations), the coefficient on \texttt{post311}
is $-0.025$ (SE $= 0.034$, $p = 0.47$), indicating no statistically or economically
significant association between 311 adoption and log per capita total general
expenditure.
A sensitivity specification restricting to cities adopting after 2012 yields a
similar estimate of $-0.020$ (SE $= 0.035$, $p = 0.56$).
Both specifications are statistically insignificant; the high table $R^2$ reflects
variance absorbed by city and year fixed effects, not the treatment indicator.

\begin{table}[htbp]
  \centering
  \caption{Effect of 311 adoption on log per capita total general expenditure.}
  \label{tab:twfe}
  \begin{threeparttable}
\begingroup
\renewcommand\cellalign{t}
\renewcommand\arraystretch{1}
\setlength{\tabcolsep}{3pt}
\begin{tabularx}{\linewidth}{@{}>{\raggedright\arraybackslash}l>{\centering\arraybackslash}X>{\centering\arraybackslash}X}
\toprule
 & \multicolumn{2}{c}{Log(exp.\ per capita)} \\
\cmidrule(lr){2-3}
 & (1) & (2) \\
\midrule
\addlinespace[1ex]
post311 & -0.025 (0.034) & -0.020 (0.035) \\
\addlinespace[0.5ex]
\midrule
\addlinespace[1ex]
 Year FE & \checkmark & \checkmark \\
\addlinespace[0.5ex]
\addlinespace[0.5ex]
City FE  & \checkmark & \checkmark \\
\addlinespace[0.5ex]
\midrule
\addlinespace[1ex]
Observations & 867 & 630 \\
\addlinespace[0.5ex]
\addlinespace[0.5ex]
$R^2$ & 0.966 & 0.95 \\
\addlinespace[0.5ex]
\bottomrule
\end{tabularx}
\endgroup
\noindent\begin{minipage}{\linewidth}\smallskip\footnotesize
Standard errors clustered at city level. Outcome: log total general expenditure per capita (Census Individual Unit Finance). Col. 1: all treated cities (adopted 1996–2022) + never-adopted controls. Col. 2: late adopters only (adopted 2012–2022) + never-adopted controls. Data: 2017–2022 only; cities adopting before 2017 are always-treated (no pre-period).\end{minipage}

\end{threeparttable}
\end{table}

% \paragraph{Event study.}
Figure~\ref{fig:eventstudy} presents event-study estimates for the subset of cities
with observable pre- and post-adoption windows in the 2017--2022 panel.
The treated group consists of 18 cities that adopted in 2019 through 2022, which are
the only cohorts with non-reference pre-period observations given the panel start
date.
The control group is the 54 never-adopted cities.
The event window is $[-3,\,+3]$ relative to the adoption year.
Pre-period coefficients at $k=-3$ and $k=-2$ are small and statistically
insignificant (0.047, SE\,$=$\,0.060; 0.012, SE\,$=$\,0.033), consistent with
parallel pre-trends.
Post-adoption coefficients are also insignificant throughout, though confidence
intervals widen as cohort support thins.
The bottom panel of the figure reports the number of treated cities contributing
to each event-time cell; the coefficient at $k=+3$ (three cities) should not be
interpreted.

\begin{figure}[H]
  \centering
  \includegraphics[width=0.78\linewidth]{figures/fig_event_study.png}
  \caption{Event-study estimates around 311 adoption.}
  \label{fig:eventstudy}
\end{figure}

% \paragraph{Interpretation and limitations.}
The null result is informative rather than the product of insufficient power alone.
311 is primarily an administrative routing and case-management system whose direct
operating costs are modest relative to total city expenditure.
Three limitations qualify the findings.
First, the panel covers 2017--2022 only; any fiscal adjustment predating 2017 is
undetectable.
Second, the event study uses 18 treated cities and is underpowered for detecting
moderate-sized effects.
Third, the regression sample skews toward larger, better-documented cities
(mean population 497k), which may differ from smaller jurisdictions in how 311
interacts with local fiscal management.

\bigskip
\noindent\textbf{AI tools.}
Claude assisted throughout: data collection used a Claude agent for structured web searches, and analysis code was developed with Claude Code. The workflow involved iterative human--AI collaboration: planning, AI-assisted execution, and author review and revision at each stage.
